% GENERATED FILE. Edit canonical lessons, then run npm run generate:books.
% canonical-source-hash: fnv1a64:97558d2e770b6830
% canonical-lessons: HI-C369-bhrastacar, HI-C369-ghus, HI-C369-samstha, HI-C369-karmacari, HI-C369-pabandi

\chapter{Corruption}
\label{ch:hi-corruption}
\hlchaptermodality{369}

\begin{chapteropening}
\textbf{By the end of this chapter:} \emph{I can talk about bribes, bans and honest institutions.}

Complete the last lesson of chapter 369: I can talk about bribes, bans and honest institutions.
\end{chapteropening}

\section[\texorpdfstring{\hi{भ्रष्टाचार} (bhraṣṭācār)}{भ्रष्टाचार (bhraṣṭācār)}]{\hi{भ्रष्टाचार} (bhraṣṭācār) — corruption}

\label{lesson:HI-C369-bhrastacar}



\begin{quote}
Before the new one: say the Hindi for discipline, then the Hindi for to polish.
\end{quote}

\subsection*{You'll want to know: \hi{भ्रष्टाचार}}
\textbf{\hi{भ्रष्टाचार}} — \emph{bhraṣṭācār} — \textquotedblleft{}corruption\textquotedblright{}.

\begin{quote}
\hi{लोग} \hi{भ्रष्टाचार} \hi{से} \hi{परेशान} \hi{हैं।}
\end{quote}

\begin{quote}
\emph{log bhraṣṭācār se pareśān haiṁ.}
\end{quote}

\begin{quote}
People are fed up with corruption.
\end{quote}

Masculine.

\begin{cousinweb}
\textbf{\hi{भ्रष्टाचार}} (\emph{bhraṣṭācār}) joins two Sanskrit words, \emph{bhraṣṭa}, \textquotedblleft{}fallen\textquotedblright{}, and \emph{ācāra}, \textquotedblleft{}conduct\textquotedblright{} (Monier-Williams): conduct that has fallen.
\end{cousinweb}

\begin{grammarlens}[title={What you've built}]
One more word for giving reasons and opinions.
\end{grammarlens}

\subsection*{Your turn}
\begin{itemize}
\raggedright
  \item \emph{Say it:} \emph{bhraṣṭācār}
  \item \emph{Say it:} \emph{bhraṣṭācār}, once more
  \item \emph{Say it:} \emph{log bhraṣṭācār se pareśān haiṁ.} — People are fed up with corruption.
  \item \emph{Recall:} say the Hindi for discipline, then the Hindi for to polish, then say \emph{bhraṣṭācār} again
\end{itemize}

\begin{tcolorbox}[breakable,colback=teal!4,colframe=teal!35!black,title={Before you move on}]
What does \hi{भ्रष्टाचार} mean? (\textquotedblleft{}Corruption\textquotedblright{}.) Say it once more.
\end{tcolorbox}
\section[\texorpdfstring{\hi{घूस} (ghūs)}{घूस (ghūs)}]{\hi{घूस} (ghūs) — a bribe (everyday)}

\label{lesson:HI-C369-ghus}



\begin{quote}
Before the new one: say the Hindi for to polish, then the Hindi for corruption.
\end{quote}

\subsection*{You'll want to know: \hi{घूस}}
\textbf{\hi{घूस}} — \emph{ghūs} — \textquotedblleft{}a bribe (everyday)\textquotedblright{}.

\begin{quote}
\hi{अधिकारी} \hi{ने} \hi{घूस} \hi{ली।}
\end{quote}

\begin{quote}
\emph{adhikārī ne ghūs lī.}
\end{quote}

\begin{quote}
The official took a bribe.
\end{quote}

Feminine.

\begin{cousinweb}
\textbf{\hi{घूस}} (\emph{ghūs}) is inherited from \emph{*ghūssā}, \textquotedblleft{}a present\textquotedblright{}, a form CDIAL reconstructs from its descendants (CDIAL 4499). CDIAL glosses the Hindi word \textquotedblleft{}bribe\textquotedblright{}; the step from a present to a bribe is not spelled out there.
\end{cousinweb}

\begin{grammarlens}[title={What you've built}]
One more word for giving reasons and opinions.
\end{grammarlens}

\subsection*{Your turn}
\begin{itemize}
\raggedright
  \item \emph{Say it:} \emph{ghūs}
  \item \emph{Say it:} \emph{ghūs}, once more
  \item \emph{Say it:} \emph{adhikārī ne ghūs lī.} — The official took a bribe.
  \item \emph{Recall:} say the Hindi for to polish, then the Hindi for corruption, then say \emph{ghūs} again
\end{itemize}

\begin{tcolorbox}[breakable,colback=teal!4,colframe=teal!35!black,title={Before you move on}]
What does \hi{घूस} mean? (\textquotedblleft{}A bribe (everyday)\textquotedblright{}.) Say it once more.
\end{tcolorbox}
\section[\texorpdfstring{\hi{संस्था} (saṁsthā)}{संस्था (saṁsthā)}]{\hi{संस्था} (saṁsthā) — an institution, an organisation}

\label{lesson:HI-C369-samstha}



\begin{quote}
Before the new one: say the Hindi for corruption, then the Hindi for a bribe (everyday).
\end{quote}

\subsection*{You'll want to know: \hi{संस्था}}
\textbf{\hi{संस्था}} — \emph{saṁsthā} — \textquotedblleft{}an institution, an organisation\textquotedblright{}.

\begin{quote}
\hi{यह} \hi{संस्था} \hi{ग़रीब} \hi{बच्चों} \hi{को} \hi{पढ़ाती} \hi{है।}
\end{quote}

\begin{quote}
\emph{yah saṁsthā ġarīb baccoṁ ko paṛhātī hai.}
\end{quote}

\begin{quote}
This organisation teaches poor children.
\end{quote}

Feminine.

\begin{cousinweb}
\textbf{\hi{संस्था}} (\emph{saṁsthā}) is a tatsama, borrowed from Sanskrit \emph{saṃsthā}, \textquotedblleft{}a standing; an establishment\textquotedblright{} (Monier-Williams).
\end{cousinweb}

\begin{grammarlens}[title={What you've built}]
One more word for giving reasons and opinions.
\end{grammarlens}

\subsection*{Your turn}
\begin{itemize}
\raggedright
  \item \emph{Say it:} \emph{saṁsthā}
  \item \emph{Say it:} \emph{saṁsthā}, once more
  \item \emph{Say it:} \emph{yah saṁsthā ġarīb baccoṁ ko paṛhātī hai.} — This organisation teaches poor children.
  \item \emph{Recall:} say the Hindi for corruption, then the Hindi for a bribe (everyday), then say \emph{saṁsthā} again
\end{itemize}

\begin{tcolorbox}[breakable,colback=teal!4,colframe=teal!35!black,title={Before you move on}]
What does \hi{संस्था} mean? (\textquotedblleft{}An institution, an organisation\textquotedblright{}.) Say it once more.
\end{tcolorbox}
\section[\texorpdfstring{\hi{कर्मचारी} (karmacārī)}{कर्मचारी (karmacārī)}]{\hi{कर्मचारी} (karmacārī) — an employee, a member of staff}

\label{lesson:HI-C369-karmacari}



\begin{quote}
Before the new one: say the Hindi for a bribe (everyday), then the Hindi for an institution.
\end{quote}

\subsection*{You'll want to know: \hi{कर्मचारी}}
\textbf{\hi{कर्मचारी}} — \emph{karmacārī} — \textquotedblleft{}an employee, a member of staff\textquotedblright{}.

\begin{quote}
\hi{बैंक} \hi{के} \hi{कर्मचारी} \hi{आज} \hi{छुट्टी} \hi{पर} \hi{हैं।}
\end{quote}

\begin{quote}
\emph{baiṅk ke karmacārī āj chuṭṭī par haiṁ.}
\end{quote}

\begin{quote}
The bank staff are on leave today.
\end{quote}

Masculine or feminine, by the person meant.

\begin{cousinweb}
\textbf{\hi{कर्मचारी}} (\emph{karmacārī}) is a tatsama, borrowed from Sanskrit \emph{karmacārin}, \textquotedblleft{}engaged in work\textquotedblright{} (Monier-Williams).
\end{cousinweb}

\begin{grammarlens}[title={What you've built}]
One more word for giving reasons and opinions.
\end{grammarlens}

\subsection*{Your turn}
\begin{itemize}
\raggedright
  \item \emph{Say it:} \emph{karmacārī}
  \item \emph{Say it:} \emph{karmacārī}, once more
  \item \emph{Say it:} \emph{baiṅk ke karmacārī āj chuṭṭī par haiṁ.} — The bank staff are on leave today.
  \item \emph{Recall:} say the Hindi for a bribe (everyday), then the Hindi for an institution, then say \emph{karmacārī} again
\end{itemize}

\begin{tcolorbox}[breakable,colback=teal!4,colframe=teal!35!black,title={Before you move on}]
What does \hi{कर्मचारी} mean? (\textquotedblleft{}An employee, a member of staff\textquotedblright{}.) Say it once more.
\end{tcolorbox}
\section[\texorpdfstring{\hi{पाबंदी} (pābandī)}{पाबंदी (pābandī)}]{\hi{पाबंदी} (pābandī) — a restriction, a ban}

\label{lesson:HI-C369-pabandi}



\begin{quote}
Before the new one: say the Hindi for an institution, then the Hindi for an employee.
\end{quote}

\subsection*{You'll want to know: \hi{पाबंदी}}
\textbf{\hi{पाबंदी}} — \emph{pābandī} — \textquotedblleft{}a restriction, a ban\textquotedblright{}.

\begin{quote}
\hi{सरकार} \hi{ने} \hi{प्लास्टिक} \hi{पर} \hi{पाबंदी} \hi{का} \hi{प्रस्ताव} \hi{रखा।}
\end{quote}

\begin{quote}
\emph{sarkār ne plāsṭik par pābandī kā prastāv rakhā.}
\end{quote}

\begin{quote}
The government proposed a ban on plastic.
\end{quote}

Feminine.

\begin{cousinweb}
\textbf{\hi{पाबंदी}} (\emph{pābandī}) is usually traced to Persian \emph{pābandī}, \textquotedblleft{}a fettering\textquotedblright{}.
\end{cousinweb}

\begin{grammarlens}[title={What you've built}]
One more word for giving reasons and opinions.
\end{grammarlens}

\subsection*{Your turn}
\begin{itemize}
\raggedright
  \item \emph{Say it:} \emph{pābandī}
  \item \emph{Say it:} \emph{pābandī}, once more
  \item \emph{Say it:} \emph{sarkār ne plāsṭik par pābandī kā prastāv rakhā.} — The government proposed a ban on plastic.
  \item \emph{Recall:} say the Hindi for an institution, then the Hindi for an employee, then say \emph{pābandī} again
\end{itemize}

\begin{tcolorbox}[breakable,colback=teal!4,colframe=teal!35!black,title={Before you move on}]
What does \hi{पाबंदी} mean? (\textquotedblleft{}A restriction, a ban\textquotedblright{}.) Say it once more.
\end{tcolorbox}
